% Partie : lambda-calcul
% Chapitre : introduction
% Section : conversion

\documentclass[../../../include/open-logic-section]{subfiles}

\begin{document}

\olfileid[fr]{lam}{int}{con}
\olsection{Conversion et réduction}

Le lambda-calcul fait intervenir trois transformations.
Une réduction est orientée. Une conversion permet de parcourir
les étapes dans les deux sens.

L'$\alpha$-conversion renomme les variables liées sans changer
les occurrences libres ni les liens entre lieurs et occurrences.
Ainsi, $\lambda x.x$ et $\lambda y.y$ représentent tous deux
la fonction identité. La \olref[lam][syn][alp]{sec} donne
les conditions précises de ce renommage. La
\olref[lam][syn][bet]{sec} présente la $\beta$-réduction.
La \olref[lam][syn][eta]{sec} présente l'$\eta$-conversion.



\end{document}
